% ===============================================================
% WEEKLY REPORT — copy this file to week02.tex, week03.tex, ...
% In Overleaf: use the file dropdown above the editor (or
% Menu > Settings > Main document) to pick which root to compile.
% ===============================================================
\documentclass[11pt,a4paper]{article}

\input{setup/project}
\input{setup/preamble}
\input{setup/weekly}

% No paragraph indentation; separate paragraphs with vertical space instead.
\setlength{\parindent}{0pt}
\setlength{\parskip}{0.6em}


\begin{document}

\weeklytitle{3}

\section{Week 3: 14--18 September 2026}

\sectionprogress

\textbf{Coronary tree graph built from the centerlines.} From the raw centerlines, points and
polyline connections per vessel side, I built a tree graph structure so future topology features could be extracted off the same structure.


\textbf{First features: angle, volume, tortuosity.} On top of that tree I started computing
bifurcation angle, angle together with local branch volume, and tortuosity, split into its several
distinct parts: the straight-line-distance ratio ($L/D$), local bending (curvature), out-of-plane
corkscrewing (torsion, non-planarity), and how often the vessel oscillates (inflections, SOAM).

\textbf{Radius attached to every centerline point.} The delivered centerlines carry no radius, so I
measure it separately at each point: the distance to the nearest background voxel in the paired
segmentation mask, i.e.\ the radius of the largest sphere centered there that still fits inside the
lumen. Both volume and bifurcation angle need this, and not just as extra context: it also sets the
window skipped past each junction before the angle itself is measured.

\textbf{Tortuosity's more interesting measures may not be trustworthy.} Curvature needs a second derivative of the centerline and torsion a third; derivatives amplify noise, so these two measures may be far less reliable than the others but they may also carry more signal.


\textbf{CAS-Net pipeline reproduced end to end.} I went through the whole pipeline and wrote up
each stage, trained from scratch for 5 epochs, then ran the full inference and evaluation off the
delivered pretrained weights and reproduced the expected results.

\textbf{Extending to Herlev--{\O}sterbro.} Herlev--{\O}sterbro comes with no delivered centerlines
and no segment names, so my plan is to extract centerlines from the 160 ImageCAS-X test
predictions and measure how far the features drift compared to the corrected masks, as a baseline
error estimate. Herlev--{\O}sterbro will be out of distribution for CAS-Net, though, so I don't yet
know how much that baseline will actually tell me about it.

\subsection{Next week}

I will spend the full week at an internal hackathon hosted by Microsoft, so sadly I won't be able
to join next week's meeting :(

\begin{itemize}
    \item Continue investigating how to extract features, now that I have predictions.
    \item Compute centerlines from the predicted segmentations and check how far the feature
    extraction deviates from the values computed on the delivered centerlines.
    \item Start looking into representation-learning approaches for the tortuosity features.
\end{itemize}

\subsection{Questions}

\begin{itemize}
    \item For building the tree and the first features, is it correct to start from the delivered
    ground-truth centerlines, before centerlines recomputed on my own predicted segmentations are
    available?
    \item For representation learning on the tree topology, patients do not share a fixed tree
    structure (branch counts and even which named vessels are present vary), so any fixed-length
    input already assumes a correspondence across patients that has not been established. Is
    solving that correspondence by hand (e.g.\ per named vessel, as planned) the right approach, or
    should this push toward a model that consumes the tree structure directly?
    \item I have not read up on this properly yet, so these are initial thoughts, tell me if this is
    not the right direction.

    If I do not do representation learning, I am bounded by whatever finite set of descriptors I
    compute on the topology; representation learning does not have that limitation. Tortuosity is, I
    think, the most important of these features, so I want to reason from that starting point
    specifically. For representation learning I would still start from the centerlines, and for
    tortuosity specifically I do not think I lose anything by working from there rather than the
    full volume.

    Concretely, would I feed the 3D point sequence of a centerline directly into an encoder, and
    check afterward whether cases like a single tight loop and a curve with many small wiggles end
    up far apart in the embedding space, even though a plain distance/chord ratio would score them
    as nearly identical? Would this be done per coronary artery, or over the whole tree at once? And
    do I need to separate by coronary dominance, so that any separation in the embedding comes from
    intrinsic tortuosity rather than from right- versus left-dominant anatomy?

    On the outcome side, if the latent space does separate different tortuosity types and only some
    of them turn out to predict the outcome, how would I then identify which geometric property is
    responsible for that difference?

    Right now the outcome can just be diseased vs.\ healthy in ImageCAS, but it would be really
    interesting if this is the direction to take Herlev--{\O}sterbro: first scan and later outcome.
    So, to see if there is a latent representation of tortuosity that can predict future plaque
    occurrence.
\end{itemize}

\include{content/methods/topology/coronary_tree_graph}
\include{content/methods/pipeline/data_preparation}
\include{content/methods/pipeline/model_architecture}
\include{content/methods/pipeline/training}
\include{content/methods/pipeline/inference}
\include{content/methods/pipeline/evaluation_metrics}

\end{document}